% Source: https://artofproblemsolving.com/community/c3h1450456_is_there_any_way_to_simplify_this
% Stable identity: aops-topic-1450456-q01
\textbf{(a)} $k!\dbinom{n}{k}=\dfrac{n!}{(n-k)!}=n(n-1)\cdots(n-k+1)$, the number of ordered selections of $k$ objects from $n$. Every such product with $k\ge2$ is $n$ times the corresponding product for $n-1$ with $k-1$ factors:
\[
S_n=n+\sum_{k=2}^{n}n\cdot(n-1)\cdots(n-k+1)=n+n\,S_{n-1}=n\,(S_{n-1}+1).
\]

\textbf{(b)} Put $j=n-k$:
\[
S_n=\sum_{k=1}^{n}\dfrac{n!}{(n-k)!}=\sum_{j=0}^{n-1}\dfrac{n!}{j!}.
\]
Now
\[
e\cdot n!=\sum_{j=0}^{n}\dfrac{n!}{j!}+\sum_{j>n}\dfrac{n!}{j!}=S_n+1+\varepsilon_n ,
\]
where the first sum is an integer ($j=n$ contributes the $1$) and
\[
0<\varepsilon_n=\dfrac{1}{n+1}+\dfrac{1}{(n+1)(n+2)}+\dots<\dfrac{1}{n+1}+\dfrac{1}{(n+1)^2}+\dots=\dfrac{1}{n}\le1 .
\]
So $S_n+1$ is the integer part of $e\cdot n!$, that is,
\[
\boxed{S_n=\lfloor e\cdot n!\rfloor-1}.
\]
For example $S_1=1$, $S_2=4$, $S_3=15$, $S_4=64$, and $\lfloor e\cdot24\rfloor=\lfloor65.2\ldots\rfloor=65$.
